\scanpage{31}
\unsure{tique} est la plus évidente. Car, si $U$ est étale sur $X$
\[
\operatorname{Fr}_{X*}(F)(U)=F(\operatorname{Fr}_X^*(U))=F(U^{(p/X)})
\]
\deleted{Or $\operatorname{Fr}_X^*(U)\simeq U^{(p/X)}$}
Or on a un \uline{isom} canonique fonctoriel en $U$
\diagram{\begin{tikzcd}U \arrow[r,"\operatorname{Fr}_{U/X}","\sim"'] & U^{(p/X)}\end{tikzcd}}
\margin{\diagram{\begin{tikzcd}[column sep=small,row sep=small]X \arrow[d,no head,"\operatorname{Fr}_X"'] & U \arrow[d,no head,"\operatorname{Fr}_U"]\\X & U \arrow[l]\end{tikzcd}}}
qui donne un isom fonctoriel en $U$
\diagram{\begin{tikzcd}\operatorname{Fr}_{X*}(F)(U)\deleted{\unsure{\simeq}}=F(U^{(p/X)}) \arrow[r,"\sim"] & F(U)\end{tikzcd}}
-- qui est l'hom cherché
\[
\boxed{\operatorname{Fr}_{F/X*}:\quad\operatorname{Fr}_{X*}(F)\longrightarrow F}
\]

\deleted{Supposons que $F$ soit représenté par $X'$ étale sur $X$

Considérons l'hom canonique

\deleted{Alors}
\diagram{\begin{tikzcd}
\operatorname{Fr}_{F/X}^*:\quad\deleted{\unsure{\operatorname{Fr}_X}}F \arrow[r] & \operatorname{Fr}_X^*F\\
\operatorname{Fr}_X^*\operatorname{Fr}_{X*}(F) \arrow[u] \arrow[ur] &
\end{tikzcd}}
Si $U$ est un \deleted{\unsure{schéma}} étale sur $X$, on définit
\diagram{\begin{tikzcd}F(U) \arrow[r] & \operatorname{Fr}_X^*(F)(U)\end{tikzcd}}
en \unsure{écrivant} $U$ sous forme de $\operatorname{Fr}_X^*(V)$, $V$ étale sur $X$, et \unsure{aura}
\diagram{\begin{tikzcd}\operatorname{Fr}_X^*\operatorname{Fr}_{X*}(F)(U) \arrow[r,leftrightarrow,"\sim"] & \operatorname{Fr}_{X*}(F)(V)\end{tikzcd}}
}

\uline{Prop} 1) $\operatorname{Fr}_{F/X}^*:\quad\operatorname{Fr}_X^*(F)\longleftrightarrow F$ pour $X,F$ variables est compatible avec extension de la base quelconque $X'\xrightarrow{g}X$, i.e. commutativité de
\margin{\diagram{\begin{tikzcd}[column sep=small,row sep=small]X \arrow[d,"\operatorname{Fr}_X"'] & X' \arrow[l,"g"'] \arrow[d,"\operatorname{Fr}_{X'}"]\\X & X' \arrow[l,"g"']\end{tikzcd}}}
\diagram{\begin{tikzcd}[column sep=huge,row sep=large]
g^*(\operatorname{Fr}_X^*(F)) & g^*(F)\overset{\mathrm{dfn}}=F' \arrow[l,"g^*(\operatorname{Fr}_{F/X}^*)"']\\
\operatorname{Fr}_{X'}^*(g^*(F))=\operatorname{Fr}_{X'}^*(F') \arrow[u,no head,"\simeq\ \mathrm{trans.}"'] \arrow[ur,leftrightarrow,"\operatorname{Fr}_{F'/X'}^*"'] &
\end{tikzcd}}
\mspage{30}
